\section{Measurement Details}
\label{sec:measurement-details}

\subsection{Timing and statistical units}

The primary MetroPT-3 endpoint is the maximum pipeline time across four independent GPU workers. Each of three execution orders is summarized by the median of two timing repeats. The reported $1.745535\times$ is the descriptive geometric mean of the three order ratios. Both paths use the same input, initialization, regularization, tolerance, verification schedule, backend, and consumer. Each repeat invokes and evaluates 16 verifier candidates and launches 32 plan-application kernels.

The two-host MetroPT-3 comparison uses three order-balanced plans on each host, with four GPU placements per host. Its six host/plan combinations are paired timing units for a fixed workload; placements and repeats stabilize timing. The $E_2$ ratio is $1.744723\times$, with stratified bootstrap interval $[1.743410,1.746706]$. Subtracting each unit's paired consumer time gives the $E_1$ ratio \CMetroRetirementOTSpeedup{}. The primary four-worker and two-host comparisons use their own endpoints and observations.

GPU energy is the change in the A100 NVML cumulative energy counter between the start of the method call and the final CUDA synchronization. The counter reports millijoules, converted to joules. This interval starts with prepared arrays and includes solver, initialization, and consumer work; compilation, warm-up, and prepared-input transfer occur before it. Unsupported counters are recorded as unavailable.

\subsection{Training and output measurements}

ImageNet-32 uses paired seeds 20260891, 20260892, and 20260893 throughout training and NFE evaluation. \Cref{tab:imagenet-solver-pairs} reports each seed's solver time. The full-training ratio uses the complete 50,000-update interval. The separate coupling-preparation diagnostic gives paired ratios $1.156$, $1.082$, and $1.079$, with fixed/DC update counts $64/52$, $64/56$, and $64/56$.

\begin{table}[!htbp]
\centering
\caption{ImageNet-32 solver times by training seed. The geometric mean of paired fixed/compacted ratios is $1.054$.}
\label{tab:imagenet-solver-pairs}
\small
\setlength{\tabcolsep}{5pt}
\renewcommand{\arraystretch}{1.15}
\begin{tabular}{@{}lrrr@{}}
\toprule
Measurement & Seed 20260891 & Seed 20260892 & Seed 20260893 \\
\midrule
Fixed time (s) & 123.429 & 124.233 & 118.840 \\
Compacted time (s) & 116.699 & 115.934 & 115.036 \\
Fixed / compacted & 1.058 & 1.072 & 1.033 \\
\bottomrule
\end{tabular}
\end{table}


All 48 worker-level primary MetroPT-3 repeats converge. The solver receives sensor inputs without failure labels. Maximum paired differences in transport cost and barycentric shift are $0.02792263$ and $0.00346756$ under the fixed failure-window consumer.

The Packer19 tolerance study uses two hosts, four GPUs per host, three method orders, and three timing repeats. Each tolerance has 24 paired observations indexed by host, order, and GPU, with an equal-host bootstrap interval. Consumer TV compares outputs to the $10^{-6}$ reference. Absolute OT-stage medians at $10^{-3}$ are 2.038 s for logical masking and 1.299 s for compaction in this stopping study. The component study's phase medians in \cref{tab:phase-costs} are summarized independently and belong to its own LM/DC comparison.

\begin{table}[!htbp]
\centering
\caption{Packer19 component-study phase times in milliseconds. Entries are separately computed phase medians. Whole-path ratios in Table~\ref{tab:packer19-components} come from paired timings.}
\label{tab:phase-costs}
\small
\setlength{\tabcolsep}{5pt}
\renewcommand{\arraystretch}{1.15}
\begin{tabular}{@{}lrrrr@{}}
\toprule
Execution & Correction & Batched verification & Verifier & Mask / compact \\
\midrule
Logical mask & 1706.3 & 638.5 & 104.5 & 0.306 \\
Compacted & 1070.8 & 400.6 & 104.5 & 3.370 \\
\bottomrule
\end{tabular}
\end{table}


\subsection{Width measurements and backend settings}

The width sweep uses an A100-SXM4-80GB with 108 SMs, driver 535.129.03, PyTorch 2.5.1 with CUDA 12.4, and Triton 3.1.0. Each primitive call contains two packed LogSumExp update kernels. Widths are randomized within 30-repeat blocks after warm-up. The $n=1024$ sweep covers widths 1--8, 10, 12, and 16; the other sizes cover 1, 2, 3, 4, 6, 8, 12, and 16.

OTT-JAX verifies at phase boundaries and rebuckets unfinished problems. The tolerance sweep uses one fixed ImageNet-32 input, ten timing repeats per tolerance, and phase sizes selected on disjoint warm-up inputs: 20 updates at $10^{-3}$ and 10 updates at the other reported tolerances. PyKeOps uses a matrix-free batched operator, geometric verification schedule, and removal of completed state. Its tiled-log implementation reaches support size $n=65536$.

The shared-initialization experiment uses a constructed stream with 90\% support overlap over real ImageNet features. Its warm replay includes descriptor, transfer, initialization, correction, and verification costs; compilation, cache loading, and process startup occur before measurement. The reported result covers 40 controlled windows.

\subsection{Worker scaling and grouping}

Independent GPU workers process different windows without a Sinkhorn collective. The one-to-four-worker throughput experiment holds work per worker fixed. The eight-worker experiment compares fixed-width and compacted execution at that worker count. \Cref{tab:additional-runtime} gives the two eight-worker inputs and the separate single-worker support-size comparisons. The reported support-size configurations meet the prespecified consumer criteria on both hosts.

\begin{table}[!htbp]
\centering
\caption{Additional MetroPT-3 runtime measurements. Eight-worker rows use fixed work per worker and pipeline time; support-size rows are separate single-worker comparisons summarized by equal-host geometric means. Energy ratios use the corresponding eight-worker endpoint.}
\label{tab:additional-runtime}
\small
\setlength{\tabcolsep}{5pt}
\renewcommand{\arraystretch}{1.15}
\begin{tabular}{@{}lrrrr@{}}
\toprule
Configuration & Fixed (s) & Compacted (s) & Time ratio & Energy ratio \\
\midrule
8 workers, higher workload & 12.5 & 4.03 & 3.11 & 3.44 \\
8 workers, lower workload & 8.00 & 3.71 & 2.16 & 2.29 \\
1 worker, $n=8192$ & --- & --- & 2.32 & --- \\
1 worker, $n=32768$ & --- & --- & 1.62 & --- \\
\midrule \multicolumn{5}{@{}l@{}}{Update counts: higher workload 4736 / 1276; lower workload 3008 / 1244.} \\
\bottomrule
\end{tabular}
\end{table}


Grouping uses a fixed multiset of 32 problems. Offline iteration counts $q_t$ determine the layouts, fixing $\sum_t q_t$ while changing $\sum_w Wq_{\max,w}$. Runtime ratios describe the resulting effect for this problem set. The layouts provide a controlled analysis of completion variation.

\subsection{Reproducibility}

The experiment artifacts record input and source identities, software and GPU environments, commands, timings, and output checks. Per-seed measurements and paired analyses are available at \url{https://github.com/cis-aconiticacid/DrainSinkhorn}. Figure sources and plotting scripts accompany the manuscript.
